% --- Archon physics formalization source begin ---
% archon:physics
% archon:covers PhyXMiniProblems/problem_phyx_mini_0174.lean
% archon:source-report reports/phyx_mini/problem_phyx_mini_0174.source.json

\chapter{Physics problem phyx\_mini\_0174}
\label{ch:PhyXMiniProblems_problem_phyx_mini_0174}

\paragraph{Problem source.}
\#\# Physical scenario

Ocean waves moving at a speed of \textbackslash{}( 4.0\textbackslash{},\textbackslash{}text\{m/s\} \textbackslash{}) are approaching a beach at angle \textbackslash{}( 	heta\_1 = 30\textasciicircum{}\textbackslash{}circ \textbackslash{}) to the normal, as shown from above in figure. Suppose the water depth changes abruptly at a certain distance from the beach and the wave speed there drops to \textbackslash{}( 3.0\textbackslash{},\textbackslash{}text\{m/s\} \textbackslash{}). Assume the same law of refraction as for light.

\#\# Question

Close to the beach, what is the angle \textbackslash{}( 	heta\_2 \textbackslash{}) between the direction of wave motion and the normal?

\#\# Answer choices

- A: A: 30.0\textasciicircum{}\textbackslash{}circ \textbackslash{}

- B: B: 41.8\textasciicircum{}\textbackslash{}circ \textbackslash{}

- C: C: 22.0\textasciicircum{}\textbackslash{}circ \textbackslash{}

- D: D: 50.0\textasciicircum{}\textbackslash{}circ \textbackslash{}

\#\# Figure caption (auxiliary; use the image as primary evidence)

The image depicts a diagram illustrating wave refraction at a boundary between deep and shallow water. 

- **Sections:**
  - **Shallow water** (light blue) at the top.
  - **Deep water** (darker blue) at the bottom.
  
- **Boundary:**
  - A dashed horizontal line separates the shallow and deep water.

- **Wave Propagation:**
  - Red lines with an arrow show wave direction transitioning from deep to shallow water.
  - The wavefronts are depicted by multiple parallel red lines perpendicular to the direction of wave propagation.

- **Angles:**
  - Two angles are labeled:
    - **θ1**: The angle between the wave direction and the vertical in deep water.
    - **θ2**: The angle between the wave direction and the vertical in shallow water.
    
- **Text Labels:**
  - **Shoreline**: Appears at the top, indicating the water's border with land.
  - **Shallow water**: Labeled in the light blue section.
  - **Deep water**: Labeled in the darker blue section.

- **Vertical Line:**
  - A dashed vertical line intersects the boundary showing angles of incidence and refraction. 

The image effectively illustrates the concept of wave refraction as waves move from deeper to shallower water, similar to Snell's Law in optics.

\#\# Dataset metadata

Wave Properties; Physical Model Grounding Reasoning, Spatial Relation Reasoning

\paragraph{Recorded answer/context.}
C: C: 22.0\textasciicircum{}\textbackslash{}circ \textbackslash{}

\paragraph{Figure/image path.}
/root/proposal\_for\_physic/science-mango/phyx\_mini\_run/phyx\_data/test\_image/174.png

\paragraph{Formalization target.}
create a compiling Lean file with sorry bodies at `PhyXMiniProblems/problem\_phyx\_mini\_0174.lean`. The Lean declarations must preserve the physical quantities, dimensions or dimensional roles, figure labels, governing-law hypotheses, and final relation expressed by this problem.
Use Mathlib/Physlib names found through LeanExplore where available. If a physics API is missing, introduce faithful local abstractions rather than scalar placeholder aliases.

\begin{theorem}[Physics formalization target]
\label{thm:physics:phyx_mini_0174:target}
\lean{PhyXMiniProblems.ProblemPhyXMini0174.problem_phyx_mini_0174}
\uses{def:physics:phyx-mini-0174:phyxminiproblems-problemphyxmini0174-oceanwaverefractionsetup, def:physics:phyx-mini-0174:phyxminiproblems-problemphyxmini0174-hasphysicalparameters, def:physics:phyx-mini-0174:phyxminiproblems-problemphyxmini0174-matchesproblemreadouts, def:physics:phyx-mini-0174:phyxminiproblems-problemphyxmini0174-matchessuppliedfigure, def:physics:phyx-mini-0174:phyxminiproblems-problemphyxmini0174-obeyswavesnelllaw, def:physics:phyx-mini-0174:phyxminiproblems-problemphyxmini0174-isuniquematchingdisplayedanswer}
The assigned autoformalize agent should translate this physics problem into Lean declarations in the covered file, with theorem and lemma proof bodies written as `by sorry`.
\end{theorem}
\begin{proof}
This is an autoformalization task, not a proof task. Produce statements that can later be proved by the physics prover without weakening the physical model.
\end{proof}

% --- Archon declaration topology begin ---
\section{Declaration topology}
\begin{definition}[Water Depth]
  \label{def:physics:phyx-mini-0174:phyxminiproblems-problemphyxmini0174-waterdepth}
  \lean{PhyXMiniProblems.ProblemPhyXMini0174.WaterDepth}
  A physical water depth carrying the dimension of length
\end{definition}
\begin{proof}
  This is the declared model definition of Water Depth.
\end{proof}
\begin{definition}[depth In Meters]
  \label{def:physics:phyx-mini-0174:phyxminiproblems-problemphyxmini0174-depthinmeters}
  \lean{PhyXMiniProblems.ProblemPhyXMini0174.depthInMeters}
  \uses{def:physics:phyx-mini-0174:phyxminiproblems-problemphyxmini0174-waterdepth}
  Read a physical water depth in metres
\end{definition}
\begin{proof}
  This is the declared model definition of depth In Meters.
\end{proof}
\begin{definition}[speed In Meters Per Second]
  \label{def:physics:phyx-mini-0174:phyxminiproblems-problemphyxmini0174-speedinmeterspersecond}
  \lean{PhyXMiniProblems.ProblemPhyXMini0174.speedInMetersPerSecond}
  Read a physical propagation speed in metres per second
\end{definition}
\begin{proof}
  This is the declared model definition of speed In Meters Per Second.
\end{proof}
\begin{definition}[angle Of Degrees]
  \label{def:physics:phyx-mini-0174:phyxminiproblems-problemphyxmini0174-angleofdegrees}
  \lean{PhyXMiniProblems.ProblemPhyXMini0174.angleOfDegrees}
  Interpret a displayed degree value as a real angle
\end{definition}
\begin{proof}
  This is the declared model definition of angle Of Degrees.
\end{proof}
\begin{definition}[angle In Degrees]
  \label{def:physics:phyx-mini-0174:phyxminiproblems-problemphyxmini0174-angleindegrees}
  \lean{PhyXMiniProblems.ProblemPhyXMini0174.angleInDegrees}
  The principal real representative of an angle, converted to degrees
\end{definition}
\begin{proof}
  This is the declared model definition of angle In Degrees.
\end{proof}
\begin{definition}[Water Region]
  \label{def:physics:phyx-mini-0174:phyxminiproblems-problemphyxmini0174-waterregion}
  \lean{PhyXMiniProblems.ProblemPhyXMini0174.WaterRegion}
  The two constant-depth water regions separated in the figure
\end{definition}
\begin{proof}
  This is the declared model definition of Water Region.
\end{proof}
\begin{definition}[Wave Segment]
  \label{def:physics:phyx-mini-0174:phyxminiproblems-problemphyxmini0174-wavesegment}
  \lean{PhyXMiniProblems.ProblemPhyXMini0174.WaveSegment}
  The incident and transmitted portions of the wave path
\end{definition}
\begin{proof}
  This is the declared model definition of Wave Segment.
\end{proof}
\begin{definition}[Figure Boundary]
  \label{def:physics:phyx-mini-0174:phyxminiproblems-problemphyxmini0174-figureboundary}
  \lean{PhyXMiniProblems.ProblemPhyXMini0174.FigureBoundary}
  The two approximately horizontal physical lines visible in the figure
\end{definition}
\begin{proof}
  This is the declared model definition of Figure Boundary.
\end{proof}
\begin{definition}[Depth Transition Model]
  \label{def:physics:phyx-mini-0174:phyxminiproblems-problemphyxmini0174-depthtransitionmodel}
  \lean{PhyXMiniProblems.ProblemPhyXMini0174.DepthTransitionModel}
  Idealizations for how water depth varies on crossing the named boundary
\end{definition}
\begin{proof}
  This is the declared model definition of Depth Transition Model.
\end{proof}
\begin{definition}[region]
  \label{def:physics:phyx-mini-0174:phyxminiproblems-problemphyxmini0174-wavesegment-region}
  \lean{PhyXMiniProblems.ProblemPhyXMini0174.WaveSegment.region}
  \uses{def:physics:phyx-mini-0174:phyxminiproblems-problemphyxmini0174-waterregion, def:physics:phyx-mini-0174:phyxminiproblems-problemphyxmini0174-wavesegment}
  The homogeneous water region occupied by each depicted wave segment
\end{definition}
\begin{proof}
  This is the declared model definition of region.
\end{proof}
\begin{definition}[Diagram Plane]
  \label{def:physics:phyx-mini-0174:phyxminiproblems-problemphyxmini0174-diagramplane}
  \lean{PhyXMiniProblems.ProblemPhyXMini0174.DiagramPlane}
  The two-dimensional Euclidean plane of the overhead diagram
\end{definition}
\begin{proof}
  This is the declared model definition of Diagram Plane.
\end{proof}
\begin{definition}[Diagram Direction]
  \label{def:physics:phyx-mini-0174:phyxminiproblems-problemphyxmini0174-diagramdirection}
  \lean{PhyXMiniProblems.ProblemPhyXMini0174.DiagramDirection}
  \uses{def:physics:phyx-mini-0174:phyxminiproblems-problemphyxmini0174-diagramplane}
  A geometric direction is represented by a nonzero diagram-plane vector
\end{definition}
\begin{proof}
  This is the declared model definition of Diagram Direction.
\end{proof}
\begin{definition}[angle Between Directions]
  \label{def:physics:phyx-mini-0174:phyxminiproblems-problemphyxmini0174-anglebetweendirections}
  \lean{PhyXMiniProblems.ProblemPhyXMini0174.angleBetweenDirections}
  \uses{def:physics:phyx-mini-0174:phyxminiproblems-problemphyxmini0174-diagramdirection}
  The undirected angle between two nonzero geometric directions
\end{definition}
\begin{proof}
  This is the declared model definition of angle Between Directions.
\end{proof}
\begin{definition}[Ocean Wave Refraction Setup]
  \label{def:physics:phyx-mini-0174:phyxminiproblems-problemphyxmini0174-oceanwaverefractionsetup}
  \lean{PhyXMiniProblems.ProblemPhyXMini0174.OceanWaveRefractionSetup}
  \uses{def:physics:phyx-mini-0174:phyxminiproblems-problemphyxmini0174-waterdepth, def:physics:phyx-mini-0174:phyxminiproblems-problemphyxmini0174-waterregion, def:physics:phyx-mini-0174:phyxminiproblems-problemphyxmini0174-wavesegment, def:physics:phyx-mini-0174:phyxminiproblems-problemphyxmini0174-figureboundary, def:physics:phyx-mini-0174:phyxminiproblems-problemphyxmini0174-depthtransitionmodel, def:physics:phyx-mini-0174:phyxminiproblems-problemphyxmini0174-diagramdirection}
  The physical quantities and named geometry in the ocean-wave refraction scenario. In particular, thetaTwo is stored as unconstrained physical data; its numerical value is not built into this structure
\end{definition}
\begin{proof}
  This is the declared model definition of Ocean Wave Refraction Setup.
\end{proof}
\begin{definition}[Is Physical Acute Angle]
  \label{def:physics:phyx-mini-0174:phyxminiproblems-problemphyxmini0174-isphysicalacuteangle}
  \lean{PhyXMiniProblems.ProblemPhyXMini0174.IsPhysicalAcuteAngle}
  An angle lies on the nonnegative acute branch used by the figure
\end{definition}
\begin{proof}
  This is the declared model definition of Is Physical Acute Angle.
\end{proof}
\begin{definition}[Has Physical Parameters]
  \label{def:physics:phyx-mini-0174:phyxminiproblems-problemphyxmini0174-hasphysicalparameters}
  \lean{PhyXMiniProblems.ProblemPhyXMini0174.HasPhysicalParameters}
  \uses{def:physics:phyx-mini-0174:phyxminiproblems-problemphyxmini0174-depthinmeters, def:physics:phyx-mini-0174:phyxminiproblems-problemphyxmini0174-speedinmeterspersecond, def:physics:phyx-mini-0174:phyxminiproblems-problemphyxmini0174-oceanwaverefractionsetup, def:physics:phyx-mini-0174:phyxminiproblems-problemphyxmini0174-isphysicalacuteangle}
  Positivity, depth ordering, and the principal angular branch. These premises select the physical solution of the sine law without assigning a numerical value to thetaTwo
\end{definition}
\begin{proof}
  This is the declared model definition of Has Physical Parameters.
\end{proof}
\begin{definition}[Matches Problem Readouts]
  \label{def:physics:phyx-mini-0174:phyxminiproblems-problemphyxmini0174-matchesproblemreadouts}
  \lean{PhyXMiniProblems.ProblemPhyXMini0174.MatchesProblemReadouts}
  \uses{def:physics:phyx-mini-0174:phyxminiproblems-problemphyxmini0174-angleofdegrees, def:physics:phyx-mini-0174:phyxminiproblems-problemphyxmini0174-oceanwaverefractionsetup}
  Numerical and qualitative information stated in the problem: an abrupt depth change, a 4.0 m/s incident wave, a 3.0 m/s transmitted wave, and theta\_1 = 30 degrees. No value of thetaTwo occurs here
\end{definition}
\begin{proof}
  This is the declared model definition of Matches Problem Readouts.
\end{proof}
\begin{definition}[Matches Supplied Figure]
  \label{def:physics:phyx-mini-0174:phyxminiproblems-problemphyxmini0174-matchessuppliedfigure}
  \lean{PhyXMiniProblems.ProblemPhyXMini0174.MatchesSuppliedFigure}
  \uses{def:physics:phyx-mini-0174:phyxminiproblems-problemphyxmini0174-angleofdegrees, def:physics:phyx-mini-0174:phyxminiproblems-problemphyxmini0174-anglebetweendirections, def:physics:phyx-mini-0174:phyxminiproblems-problemphyxmini0174-oceanwaverefractionsetup}
  Geometric information read from the primary image. The depth boundary and shoreline are parallel, the dashed line is normal to the depth boundary, each drawn wavefront is perpendicular to its propagation direction, and the two theta labels are the corresponding angles from that same normal. These relations interpret thetaTwo geometrically but do not prescribe its value
\end{definition}
\begin{proof}
  This is the declared model definition of Matches Supplied Figure.
\end{proof}
\begin{definition}[Obeys Wave Snell Law]
  \label{def:physics:phyx-mini-0174:phyxminiproblems-problemphyxmini0174-obeyswavesnelllaw}
  \lean{PhyXMiniProblems.ProblemPhyXMini0174.ObeysWaveSnellLaw}
  \uses{def:physics:phyx-mini-0174:phyxminiproblems-problemphyxmini0174-speedinmeterspersecond, def:physics:phyx-mini-0174:phyxminiproblems-problemphyxmini0174-oceanwaverefractionsetup}
  The wave form of Snell's law at a depth boundary, sin(theta\_1) / v\_1 = sin(theta\_2) / v\_2, written without division as sin(theta\_1) * v\_2 = sin(theta\_2) * v\_1 using SI speed readouts. Mathlib and Physlib provide the angle, sine, and dimensionful speed objects, but no declaration specialized to this water-wave refraction law
\end{definition}
\begin{proof}
  This is the declared model definition of Obeys Wave Snell Law.
\end{proof}
\begin{lemma}[theta Two sine eq three eighths]
  \label{lem:physics:phyx-mini-0174:phyxminiproblems-problemphyxmini0174-thetatwo-sine-eq-three-eighths}
  \lean{PhyXMiniProblems.ProblemPhyXMini0174.thetaTwo_sine_eq_three_eighths}
  \uses{def:physics:phyx-mini-0174:phyxminiproblems-problemphyxmini0174-oceanwaverefractionsetup, def:physics:phyx-mini-0174:phyxminiproblems-problemphyxmini0174-matchesproblemreadouts, def:physics:phyx-mini-0174:phyxminiproblems-problemphyxmini0174-obeyswavesnelllaw}
  The given speeds and 30 degree incident angle reduce the wave Snell law to sin(theta\_2) = 3/8
\end{lemma}
\begin{proof}
  The conclusion follows from the governing relations and figure readouts named in the statement of theta Two sine eq three eighths.
\end{proof}
\begin{definition}[Answer Choice]
  \label{def:physics:phyx-mini-0174:phyxminiproblems-problemphyxmini0174-answerchoice}
  \lean{PhyXMiniProblems.ProblemPhyXMini0174.AnswerChoice}
  Labels of the four answers displayed with the problem
\end{definition}
\begin{proof}
  This is the declared model definition of Answer Choice.
\end{proof}
\begin{definition}[direction Degrees]
  \label{def:physics:phyx-mini-0174:phyxminiproblems-problemphyxmini0174-answerchoice-directiondegrees}
  \lean{PhyXMiniProblems.ProblemPhyXMini0174.AnswerChoice.directionDegrees}
  \uses{def:physics:phyx-mini-0174:phyxminiproblems-problemphyxmini0174-answerchoice}
  The angle in degrees printed beside each displayed answer label
\end{definition}
\begin{proof}
  This is the declared model definition of direction Degrees.
\end{proof}
\begin{definition}[recorded Answer Choice]
  \label{def:physics:phyx-mini-0174:phyxminiproblems-problemphyxmini0174-recordedanswerchoice}
  \lean{PhyXMiniProblems.ProblemPhyXMini0174.recordedAnswerChoice}
  \uses{def:physics:phyx-mini-0174:phyxminiproblems-problemphyxmini0174-answerchoice}
  The answer label recorded in the supplied dataset metadata
\end{definition}
\begin{proof}
  This is the declared model definition of recorded Answer Choice.
\end{proof}
\begin{definition}[Matches Displayed Answer]
  \label{def:physics:phyx-mini-0174:phyxminiproblems-problemphyxmini0174-matchesdisplayedanswer}
  \lean{PhyXMiniProblems.ProblemPhyXMini0174.MatchesDisplayedAnswer}
  \uses{def:physics:phyx-mini-0174:phyxminiproblems-problemphyxmini0174-angleindegrees, def:physics:phyx-mini-0174:phyxminiproblems-problemphyxmini0174-answerchoice}
  Agreement with a displayed angle rounded to the nearest tenth of a degree
\end{definition}
\begin{proof}
  This is the declared model definition of Matches Displayed Answer.
\end{proof}
\begin{definition}[Is Unique Matching Displayed Answer]
  \label{def:physics:phyx-mini-0174:phyxminiproblems-problemphyxmini0174-isuniquematchingdisplayedanswer}
  \lean{PhyXMiniProblems.ProblemPhyXMini0174.IsUniqueMatchingDisplayedAnswer}
  \uses{def:physics:phyx-mini-0174:phyxminiproblems-problemphyxmini0174-answerchoice, def:physics:phyx-mini-0174:phyxminiproblems-problemphyxmini0174-matchesdisplayedanswer}
  A displayed answer is the unique candidate within the rounding tolerance
\end{definition}
\begin{proof}
  This is the declared model definition of Is Unique Matching Displayed Answer.
\end{proof}
% --- Archon declaration topology end ---
% --- Archon physics formalization source end ---
